Due to the inherent variance in the BLE protocol and high interference demonstrated on wireless channels commonly used in the urban settings used for the empirical experiment, some RSSI values should be collected again, as visible from jittery subplots in the PDF plot~\ref{fig:rssi-pdf}: for example, the 60m and 90m plot exhibit repeating jumps in their density values, in comparison to the smoothness in lower ranges (e.g. 10, 20, 30m ranges).
This we attribute to the low number of received packets at those distances (sometimes below 100 packets~\ref{fig:advertisement-success-probability}), which makes the RSSI distribution more sensitive to outliers.
The RSSI values we have provided for the ranges that are complete and stable, are mostly sufficient to derive the necessary path loss parameters for the simulation, as they exhibit the expected power-law distribution and are mostly consistent. 
% A point for future research could be to measure the maximum achieved distance with varying transmitting power levels using a different chip or a wireless dongle.
Furthermore, we were not able to have fine-grain control over exactly when the advertisements were sent, because the BlueZ Bluetooth communication stack and the underlying hardware controller handle the advertisement scheduling internally to prevent antenna overheating.
To best mitigate this, we set the advertisement interval window to the minimum possible values of to achieve the highest consistent frequency of advertisements sent.
We were however able to prove that the advertisements were not actually being consistently received at the configured interval window of 20ms using a Bluetooth sniffer on a mobile phone, which measured the advertisement intervals.
Therefore, another prospect for this project in the future, would be to use a custom-made BLE firmware on a microcontroller platform, such as the Nordic nRF52 series, to have more control over the advertisement scheduling and transmission power.
This would additionally allow us to measure the effect setting the transmission power has on the signal propagation distance, which can allow us to model the resilience of p2p networks built using different transmission power levels.